\chapter{Taxonomie des intentions}
\label{an:taxonomie}

Les 77 intentions de Banking77, rangées en neuf familles (vérifié par script : chacune une fois et
une seule), puis les huit intentions que nous créons.

\begin{center}\footnotesize
\begin{longtable}{@{}L{3.1cm} L{11.3cm}@{}}
\toprule
\textbf{Famille} & \textbf{Intentions} \\
\midrule
\endhead
\f{carte\_perdue} (4) & \f{lost\_or\_stolen\_card}, \f{compromised\_card}, \f{card\_swallowed},
  \f{lost\_or\_stolen\_phone} \\
\addlinespace
\f{carte\_usage} (10) & \f{card\_not\_working}, \f{contactless\_not\_working}, \f{pin\_blocked},
  \f{change\_pin}, \f{activate\_my\_card}, \f{card\_about\_to\_expire}, \f{card\_acceptance},
  \f{declined\_card\_payment}, \f{virtual\_card\_not\_working}, \f{pending\_card\_payment} \\
\addlinespace
\f{carte\_commande} (12) & \f{card\_arrival}, \f{card\_delivery\_estimate},
  \f{order\_physical\_card}, \f{get\_physical\_card}, \f{getting\_spare\_card},
  \f{getting\_virtual\_card}, \f{get\_disposable\_virtual\_card}, \f{disposable\_card\_limits},
  \f{card\_linking}, \f{visa\_or\_mastercard}, \f{supported\_cards\_and\_currencies},
  \f{apple\_pay\_or\_google\_pay} \\
\addlinespace
\f{frais\_contestation} (14) & \f{extra\_charge\_on\_statement}, \f{card\_payment\_fee\_charged},
  \f{transfer\_fee\_charged}, \f{cash\_withdrawal\_charge}, \f{exchange\_charge},
  \f{top\_up\_by\_card\_charge}, \f{top\_up\_by\_bank\_transfer\_charge},
  \f{transaction\_charged\_twice}, \f{card\_payment\_not\_recognised},
  \f{direct\_debit\_payment\_not\_recognised}, \f{cash\_withdrawal\_not\_recognised},
  \f{request\_refund}, \f{Refund\_not\_showing\_up}, \f{reverted\_card\_payment?} \\
\addlinespace
\f{virement} (10) & \f{transfer\_into\_account}, \f{cancel\_transfer}, \f{pending\_transfer},
  \f{transfer\_timing}, \f{failed\_transfer}, \f{declined\_transfer}, \f{beneficiary\_not\_allowed},
  \f{transfer\_not\_received\_by\_recipient}, \f{receiving\_money},
  \f{balance\_not\_updated\_after\_bank\_transfer} \\
\addlinespace
\f{compte\_identite} (8) & \f{edit\_personal\_details}, \f{verify\_my\_identity},
  \f{why\_verify\_identity}, \f{unable\_to\_verify\_identity}, \f{verify\_source\_of\_funds},
  \f{passcode\_forgotten}, \f{terminate\_account}, \f{age\_limit} \\
\addlinespace
\f{especes} (5) & \f{atm\_support}, \f{pending\_cash\_withdrawal}, \f{declined\_cash\_withdrawal},
  \f{wrong\_amount\_of\_cash\_received}, \f{balance\_not\_updated\_after\_cheque\_or\_cash\_deposit} \\
\addlinespace
\f{change\_etranger} (6) & \f{exchange\_rate}, \f{card\_payment\_wrong\_exchange\_rate},
  \f{wrong\_exchange\_rate\_for\_cash\_withdrawal}, \f{exchange\_via\_app},
  \f{fiat\_currency\_support}, \f{country\_support} \\
\addlinespace
\f{alimenter\_compte} (8) & \f{automatic\_top\_up}, \f{pending\_top\_up}, \f{top\_up\_limits},
  \f{top\_up\_reverted}, \f{topping\_up\_by\_card}, \f{top\_up\_failed}, \f{verify\_top\_up},
  \f{top\_up\_by\_cash\_or\_cheque} \\
\midrule
\f{epargne} (1) & \f{epargne\_projet} \\
\f{credit} (3) & \f{credit\_immobilier}, \f{credit\_conso}, \f{pret\_etudiant} \\
\f{virement} (+1) & \f{faire\_virement} \\
\f{aide} (1) & \f{question\_definition} \\
\f{insatisfaction} (1) & \f{insatisfaction} \\
\f{conversation} (1) & \f{hors\_sujet} \\
\bottomrule
\end{longtable}
\end{center}

Total : 14 familles, 85 intentions. Les noms \f{Refund\_not\_showing\_up} et
\f{reverted\_card\_payment?} sont repris tels qu'ils figurent dans Banking77.

\chapter{Le catalogue de produits}
\label{an:catalogue}

Produits simples, conditions publiques. \textbf{Chaque fiche du fichier JSON porte la source
publique et la date du relevé} ; les valeurs ci-dessous sont à vérifier à la source le jour de la
saisie, et les taux réglementés, qui changent, n'y figurent volontairement pas.

\begin{center}\small
\begin{longtable}{@{}L{2.5cm} L{3.1cm} L{2.6cm} L{2.2cm} L{3.7cm}@{}}
\toprule
\textbf{Produit} & \textbf{Pour qui} & \textbf{Plafond, limites} & \textbf{Argent disponible}
  & \textbf{Convient pour} \\
\midrule
\endhead
Livret A & Tout le monde, mineurs compris ; un par personne & 22\,950\eur{} & Oui & Épargne de
  précaution ; épargne d'un enfant \\
LDDS & Majeurs domiciliés fiscalement en France ; un par personne & 12\,000\eur{} & Oui &
  Compléter le Livret A \\
LEP & Revenu fiscal de référence sous un plafond annuel (à relever) & 10\,000\eur{} & Oui &
  Revenus modestes : taux plus élevé \\
Livret Jeune & 12 à 25 ans & 1\,600\eur{} & Oui & Premier livret d'un jeune \\
PEL & Tout le monde ; versements réguliers obligatoires & 61\,200\eur{} & Non : un retrait clôt le
  plan & Projet immobilier à plusieurs années \\
CEL & Tout le monde & 15\,300\eur{} & Oui & Petit projet, complément du PEL \\
\midrule
Crédit immobilier & Majeurs & Mensualités $\leq$ 35\,\% des revenus nets, assurance comprise ;
  durée $\leq$ 25 ans (27 dans certains cas) & --- & Achat d'un logement ; décision par un
  conseiller \\
Prêt à taux zéro & Primo-accédants, sous conditions de ressources et de zone & Selon conditions
  & --- & Complément d'un crédit immobilier ; éligibilité vérifiée par un conseiller \\
Prêt personnel & Majeurs & Capacité de remboursement ; mention légale obligatoire & --- & Voiture,
  travaux, équipement \\
Prêt étudiant & Étudiants majeurs & Montant et caution selon l'offre ; remboursement différé
  possible ; version garantie par l'État sous conditions & --- & Financer ses études \\
\midrule
Offre clientèle fragile & Personnes en situation de fragilité financière & Tarif mensuel plafonné ;
  commission d'intervention limitée à 4\eur{} par opération et 20\eur{} par mois & --- &
  Inclusion ; proposée par un conseiller \\
Opposition carte (service) & Tout titulaire d'une carte & --- & Immédiat & Perte, vol, fraude \\
\bottomrule
\end{longtable}
\end{center}

\paragraph{Règle générale sur les frais.} Hors offre clientèle fragile, la commission
d'intervention est plafonnée par la loi à 8\eur{} par opération et 80\eur{} par mois (article
R312-4-1 du Code monétaire et financier). C'est un plafond, pas un prix : le montant réellement
facturé est celui de la grille tarifaire de la banque, à relever sur son site.

\chapter{Le détail des parcours}
\label{an:parcours}

P : information lue dans le profil (confirmée si simulée). Q : question posée au client.
E : entité extraite du message.

\begin{center}\small
\begin{longtable}{@{}L{2.4cm} L{4.4cm} L{3.5cm} L{4.1cm}@{}}
\toprule
\textbf{Parcours} & \textbf{Informations} & \textbf{Règles, outils} & \textbf{Conclusion} \\
\midrule
\endhead
Carte perdue & Carte concernée (P ; Q s'il y en a plusieurs) ; perte ou vol (Q) ; quand (E) &
  \f{bloquer\_carte}, avec confirmation & Référence, heure ; rassurer ; proposer une nouvelle carte ;
  rappel : ne jamais communiquer son code \\
\addlinespace
Frais & Période (E ou P) ; frais visé (Q s'il y en a plusieurs) & \f{expliquer\_frais} ; plafond
  légal & Explication datée ; si les incidents se répètent, orientation vers l'offre clientèle
  fragile ; transmission si contestation \\
\addlinespace
Virement & Bénéficiaire (P ; Q) ; montant (E ; Q) ; date (Q, immédiat par défaut) &
  \f{preparer\_virement} ; solde suffisant ; bénéficiaire enregistré & Récapitulatif, confirmation,
  référence ; mode mots simples ; rappel contre l'arnaque \\
\addlinespace
Épargne & Pour qui (P : enfants ; Q) ; horizon (Q) ; montant mensuel (Q) ; épargne existante (P)
  & Éligibilité (âge, revenu, plafonds atteints) ; pertinence (horizon court : livrets ; projet
  immobilier : PEL, CEL) ; appétence & Propositions classées avec raisons ; produits écartés \\
\addlinespace
Crédit immobilier & Revenus (P) ; crédits en cours (P) ; prix (Q) ; apport (P, confirmé) & 35\,\% ;
  25 ans ; frais de notaire estimés ; capacité & Estimation, leviers, rendez-vous ; mention
  légale \\
\addlinespace
Crédit à la consommation & Objet (E) ; montant (Q) ; durée (Q) ; revenus et crédits (P) & Taux
  d'effort, crédit en cours compris ; reste à vivre (paramètre interne) & Estimation prudente ;
  mention légale ; rendez-vous \\
\addlinespace
Prêt étudiant & Statut (P) ; études et durée (Q) ; montant (Q) ; caution possible (Q) &
  Éligibilité (majeur, étudiant) & Orientation vers le prêt étudiant ; Livret Jeune ;
  rendez-vous \\
\addlinespace
Jargon & Terme (E, par le glossaire) & Recherche dans le glossaire, en mémoire & Fiche courte ;
  proposition d'un exemple \\
\addlinespace
Insatisfaction & Contacts du jour (P) & Règle des trois contacts ; ton & Excuses ; reformulation
  du problème ; \f{transmettre\_conseiller} avec résumé \\
\bottomrule
\end{longtable}
\end{center}

\chapter{Le glossaire}
\label{an:glossaire}

Fiches courtes, écrites par nous et relues ; c'est le texte que l'assistant affiche, sans LLM.

\begin{center}\small
\begin{longtable}{@{}L{3.3cm} L{11.1cm}@{}}
\toprule
\textbf{Terme} & \textbf{Explication en mots simples} \\
\midrule
\endhead
Agios & Les intérêts et frais que la banque facture quand votre compte est à découvert. \\
Apport personnel & L'argent que vous mettez vous-même dans un achat, sans l'emprunter. \\
Assurance emprunteur & Une assurance qui rembourse tout ou partie de votre crédit si vous ne
  pouvez plus le faire (décès, invalidité\dots). Elle est presque toujours demandée pour un crédit
  immobilier. \\
Commission d'intervention & Des frais facturés quand la banque doit traiter une opération qui pose
  problème, par exemple un paiement alors que votre solde est insuffisant. La loi les plafonne. \\
Découvert autorisé & La possibilité, prévue avec votre banque, d'avoir un solde négatif jusqu'à
  un certain montant. Il coûte des agios. \\
Épargne de précaution & Une réserve d'argent disponible à tout moment pour faire face à un
  imprévu ; on conseille souvent trois à six mois de dépenses. \\
Frais de notaire & Les taxes et frais payés lors de l'achat d'un logement, en plus du prix :
  environ 7 à 8\,\% dans l'ancien, 2 à 3\,\% dans le neuf. \\
IBAN & Le numéro qui identifie un compte bancaire ; il faut le donner pour recevoir un virement. \\
LDDS & Livret de développement durable et solidaire : un livret d'épargne disponible, sans impôt
  sur les intérêts, plafonné à 12\,000\eur{}, réservé aux majeurs. \\
LEP & Livret d'épargne populaire : un livret mieux rémunéré, réservé aux personnes dont les
  revenus sont sous un plafond. \\
Livret A & Le livret d'épargne le plus courant : argent disponible à tout moment, intérêts sans
  impôt, taux fixé par l'État, plafond de 22\,950\eur{}. \\
Livret Jeune & Un livret pour les 12--25 ans, plafonné à 1\,600\eur{}. \\
Opposition & Le blocage définitif d'une carte perdue ou volée, pour que personne ne puisse s'en
  servir. \\
PEL & Plan épargne logement : on verse régulièrement pendant plusieurs années, en vue d'un projet
  immobilier ; un retrait ferme le plan. \\
Prélèvement & Un paiement demandé directement à votre banque par une entreprise que vous avez
  autorisée (énergie, téléphone\dots). \\
Reste à vivre & Ce qui vous reste chaque mois une fois payés vos charges et vos crédits. \\
TAEG & Taux annuel effectif global : le coût total d'un crédit (intérêts, frais, assurance
  obligatoire) exprimé en pourcentage par an. C'est le chiffre qui permet de comparer deux offres. \\
Taux d'endettement & La part de vos revenus qui part chaque mois dans vos crédits. Pour un crédit
  immobilier, elle ne doit pas dépasser 35\,\%. \\
\bottomrule
\end{longtable}
\end{center}

\chapter{Commandes utiles}
\label{an:commandes}

\paragraph{Ollama sur un poste de la salle} (d'après les slides de Denis Barthou).

\begin{code}
\footnotesize
\begin{verbatim}
export PATH=/opt/ollama/bin:${PATH}
export OLLAMA_MODELS=/opt/ollama/models
export OLLAMA_CONTEXT_LENGTH=4096          # au plus juste (16384 par defaut sur les slides)
ollama serve &                             # un seul serveur par machine
ollama pull ministral-3:3b                 # en filaire
ollama list
\end{verbatim}
\end{code}

\paragraph{Appel à Ollama avec mesure.} La réponse contient \f{prompt\_eval\_count},
\f{eval\_count} et les durées en nanosecondes, ce qui alimente la trace.

\begin{code}
\footnotesize
\begin{verbatim}
curl -s localhost:11434/api/chat -d '{
  "model": "ministral-3:3b",
  "messages": [{"role": "system", "content": "..."},
               {"role": "user", "content": "..."}],
  "stream": false,
  "keep_alive": "30m",
  "options": {"temperature": 0.2, "num_ctx": 4096, "num_predict": 200}
}'
# Pour un modele qui raisonne (qwen3) : ajouter  "think": false
\end{verbatim}
\end{code}

\paragraph{Serveur vLLM sur DeepEirb.} vLLM expose une API compatible OpenAI ; l'adresse et le nom
du modèle sont donnés par Denis Barthou.

\begin{code}
\footnotesize
\begin{verbatim}
curl -s http://<deepeirb>:<port>/v1/chat/completions \
  -H "Content-Type: application/json" \
  -d '{"model": "<modele>", "temperature": 0,
       "messages": [{"role": "user", "content": "Traduis en francais ..."}]}'
\end{verbatim}
\end{code}

\paragraph{Lancer l'interface et les mesures.} Depuis la racine du dépôt, une fois les modules
en place :

\begin{code}
\footnotesize
\begin{verbatim}
python3 -m venv .venv && source .venv/bin/activate
pip install sentence-transformers scikit-learn streamlit pyyaml codecarbon pandas
streamlit run app.py
python eval/run_scenarios.py          # rejoue les 30 scenarios
python eval/bench_nlu.py              # tableau de comprehension
python eval/bench_frugalite.py A B C  # comparaison des configurations
\end{verbatim}
\end{code}
